\section{Hardware und Rendering}
\subsection{Verwendete Hardware}\label{Hardware}
\begin{table}[h]
    \centering
\begin{tabular}{@{}lccc@{}}
    \toprule
    \textbf{Produkt}         & \textbf{ThinkPad P15 Gen 2}                       & \textbf{ThinkPad P14s Gen 1}           & \textbf{ASUS ZenBook 13}            \\
    \midrule
    \textbf{Prozessor}       & \makecell{Intel Core \\ i7-11850H}          & \makecell{AMD Ryzen 7 \\ PRO 4750U}    & \makecell{AMD Ryzen 7 \\ 5800 U}  
    \vspace{2mm} \\ % Abstand zwischen Sektionen
    \textbf{Taktung (in GHz)}&                                             &                                        &                                     \\ 
    minimale Frequenz:       & 2.5                                         & 1.7                                    & 1.9                                 \\ 
    Boost (ein Kern):        & 4.8                                         & 4.1                                    & 4.4                                 \\ 
    Boost (alle Kerne):      & (keine Angabe)                              & 3.1                                    & 3.4                       
    \vspace{3mm} \\ % Abstand zwischen Sektionen 
    \textbf{Parallelität}    &                                             &                                        &                                     \\
    Kerne                    & 8                                           & 8                                      & 8                                   \\
    Threads                  & 16                                          & 16                                     & 16 
    \vspace{3mm} \\ % Abstand zwischen Sektionen
    \textbf{RAM}             & 32 GB                                       & 24 GiB                                 & 16 GB                               \\ 
    Taktung:                 & 3200 MHz                                    & 3200 MHz                               & 3200 MHz                   
    \vspace{3mm}  \\ % Abstand zwischen Sektionen
    \textbf{Grafikkarte}     & \makecell{Nvidia RTX \\ A3000}              & \makecell{AMD Radeon \\ RX Vega 7}     & \makecell{AMD Radeon \\ RX Vega 8}  \\ 
    Größe:                   & 6 GB                                        & 2 GB                                   & 497 MB                     
    \vspace{3mm} \\ % Abstand zwischen Sektionen
    \textbf{Betriebssystem}  & Windows 10                                  & \makecell{Debian 13 \\ Kernel 6.12.63} & Windows 10                          \\ 
    \bottomrule
\end{tabular}
    \caption{Überblick der Gerätespezifikationen verwendeter Laptops}
\end{table}
\noindent
Während der Prozessor-Boost mit einem Kern für die Single-Thread-Performance relevant ist, ist der Boost mit mehreren Kernen 
für die Multi-Thread-Performance von Bedeutung.
\vspace{0.25cm}\\
\noindent
\textit{Anmerkung:} Auf dem Thinkpad P14s Gen 1 wurde OpenCL nicht nativ auf der Grafikkarte, sondern über PoCL\footnote{PoCL = Portable Computing Language --- 
u.a. für x86-Architektur geeignete Implementierung von OpenCL} auf der CPU ausgeführt. Grund dafür waren Kooperationsprobleme zwischen der integrierten 
AMD-Grafikkarte und den Linux-Treibern, die zu Systemabstürzen führten. Da dieses Problem ausschließlich auf dem Thinkpad P14s Gen 1 unter Linux auftrat und nicht
bei Nvidia unter Linux oder ähnlichen AMD-Grafikchips unter Windows, deutet dies auf ein Treiberproblem hin, nicht auf ein Problem mit unserem Programm.

\subsection{Crossplatform mit Linux und Windows}
Das Projekt selbst ist lauffähig unter Windows und Linux, sofern \lstinline{nvidia-opencl-dev} (o.ä. für die jeweilig verbaute Hardware) 
und \lstinline{libsdl2-dev} installiert sind. \\
Im Makefile werden die Bibliotheken wie folgt verwendet: \lstinline{LFLAGS = -lSDL2 -lSDL2main -lOpenCL} 
Für die Verwendung von OpenMP wird zusätzlich \lstinline{-fopenmp} bei den CXXFLAGS benötigt --- das Makefile des Linux-Projekts 
selbst stammt aus einem Plugin für die Projektgenerierung für C++.  
Da die Abgabe unter Windows mit Visual Studio 2022 erfolgen soll, haben wir den src-Ordner vom Visual Studio Projekt ins 
Linux-Projektverzeichnis verlinkt. \\
Eine Herausforderung während der Entwicklung für beide Plattformen war vor allem die Zeitmessung. Unter Linux wurde bei 
paralleler Ausführung anfangs die Zeit für alle Threads aufsummiert, was dann quasi der Zeit entsprach, die die serielle Version
benötigte. Unter Windows hingegen wurde die Zeit korrekt gemessen, allerdings fehlten hier die Funktionen, die zum Messen der 
Zeit unter Linux nötig waren.\\
Unsere Lösung war, die jeweilige Zeitmessung in einem Präprozessor-Block zu kapseln, sodass jede Plattform ihre eigene
Implementierung verwenden kann. 

\subsection{Renderer}
Ein Pythagorasbaum definiert sich durch sein Aussehen. Dementsprechend benötigten wir zur Vollendung unseres Projekts eine Möglichkeit,
unsere Ergebnisse visuell aufzubereiten. Da das Rendern von Grafiken kein wesentlicher Bestandteil der Lehrveranstaltung war, ließen wir uns
bei der Entwicklung eines Renderers durch Chat-GPT 5 unterstützen. \\
Unser Renderer verwendet zur Darstellunug der Daten die SDL-Bibliothek\footnote{SDL = Simple DirectMedia Layer; freie Multimedia-Bibliothek mit 
Unterstützung vieler Plattformen}.\\
\noindent
Die Punkte werden dort mittels einer Projektionsmatrix auf die Bildfläche übertragen. Die so entstehenden Bildpunkte 
werden anschließend mit der Information über die Verbindung untereinander (Kanten, engl. `Edges`) im dreidimensionalen 
Raum zu einer Wireframe-Grafik\footnote{Wireframe = Grafik, welche ausschließlich aus Kanten/Ecken besteht und keine 
Flächen beinhaltet} verarbeitet, da jede Verbindung im 3D-Raum direkt auf die 2D-Bildfläche übertragen werden kann. 
So entsteht eine 2D-Projektion, welche auf dem Bildschirm angezeigt wird. Die Kamera kann mit Hilfe von WASD 
nach links, rechts, oben und unten bewegt werden, wobei die Kamera immer auf einen Fixpunkt im Koordinatensystem
gerichtet bleibt.\\
Eine Redundanzprüfung, welche verhindern würde, dass übereinanderliegende und dadurch nicht sichtbare Punkte und 
Linien gerendert werden, haben wir --- obwohl sie die Effizienz des Renderers maßgeblich steigern könnte --- nicht 
implementiert, um uns auf die tatsächlichen Anforderungen des Kurses zu fokussieren. Um ein besseres visuelles 
Erlebnis zu ermöglichen, haben wir die tatsächlich aufeinanderliegenden Kanten mit der Funktion 
\texttt{removeRedundantEdges()} entfernt. 

% - mit Hilfe von KI erstellt
% - erstellt mit SDL (Simple Direct Media Layer) eine Wireframe-Grafik (nur Kanten/Ecken, keine Flächen)
% - verarbeitet Punkte und Informationen über Verbindung selbiger (Edges) im dreidimensionalen Raum zu einer Projektion in 2D 
% - auch doppelte Punkte / Linien werden ohne Prüfung der Sichtbarkeit verarbeitet -> keine Effizienzsteigerung durch Redundanzprüfung

% Einbindung ins Projekt? 

